\documentclass[11pt]{article}

% Change "review" to "final" to generate the final version.
\usepackage[review]{acl}

% Standard package includes
\usepackage{times}
\usepackage{latexsym}
\usepackage[T1]{fontenc}
\usepackage[utf8]{inputenc}
\usepackage{microtype}
\usepackage{inconsolata}
\usepackage{graphicx}
\usepackage{booktabs}
\usepackage{amsmath}
\usepackage{amssymb}
\usepackage{multirow}
\usepackage{subcaption}

\title{Who Wrote This? Identity Bias and Context Anchoring in LLM Self-Critique}

\author{Kim Jun-yeop \\
  Affiliation \\
  \texttt{email@domain} \\}

\begin{document}
\maketitle

\begin{abstract}
Large language models (LLMs) often fail to correct their own mistakes, a phenomenon frequently attributed to inherent reasoning limitations. However, we hypothesize that this "self-correction blind spot" is primarily driven by \emph{identity bias} and \emph{context anchoring} rather than a lack of logical capacity. We introduce the \textbf{Identity Swap Critique} framework to systematically isolate these factors by presenting identical erroneous solutions under varying authorship labels (e.g., Self, External Model, Human). Our experiments across diverse reasoning benchmarks (GSM8K, MATH, GPQA) reveal a significant performance gap: models consistently identify and rectify the same logical fallacies when attributed to others, yet overlook them when labeled as their own. Furthermore, we decompose these biases through context-separation experiments, demonstrating that simple protocol-level interventions—such as anonymization and session decoupling—can recover lost reasoning performance. Our findings reframe the self-correction challenge as a protocol design problem, providing a new perspective on building reliable autonomous reasoning systems.
\end{abstract}

\section{Introduction}
\label{sec:intro}

The ability to self-correct is a hallmark of advanced intelligence. While Large Language Models (LLMs) have shown remarkable progress in complex reasoning, their capacity for autonomous error correction remains surprisingly fragile. Recent studies suggest that intrinsic self-correction often yields marginal improvements or even degrades performance without external feedback \cite{huang2024self, kamoi2024survey}. This has led to a prevailing view that models lack the "meta-cognitive" depth required to scrutinize their own reasoning paths.

In this paper, we challenge this assumption. We argue that the failure of self-correction is not a fundamental reasoning deficit, but a byproduct of \textbf{Identity Bias} and \textbf{Context Anchoring}. Specifically, we observe that models exhibit a "blind spot" when reviewing their own outputs—a tendency to rationalize existing errors rather than critically evaluate them. This mirrors human cognitive biases, such as the \emph{myside bias}, where individuals evaluate arguments more favorably when they align with their own identity or prior commitments.

To empirically investigate this, we propose the \textbf{Identity Swap Critique} framework. By maintaining the content of a flawed reasoning chain while manipulating the attributed author (e.g., "This is your solution" vs. "This is a student's solution"), we isolate the impact of social and identity framing on LLM evaluation. 

Our contributions are as follows:
\begin{itemize}
    \item We introduce the \textbf{Identity Swap Critique} framework to quantify the impact of authorship attribution on model critique performance.
    \item We provide a \textbf{bias decomposition} analysis, distinguishing between the effects of identity-driven preference and session-based context anchoring.
    \item We demonstrate that \textbf{protocol-level mitigations}, such as anonymization and session separation, significantly close the self-correction gap without requiring additional training or external verifiers.
    \item We provide a new theoretical framing: self-correction is a \textbf{perspective problem}, not just a reasoning problem.
\end{itemize}

\section{Related Work}
\label{sec:related}

\paragraph{Limits of Self-Correction} 
Prior work has documented the "self-correction blind spot" \cite{tsui2025self}, noting that models struggle to revise their own outputs. \citet{huang2024self} showed that without external labels, models often flip correct answers to incorrect ones during self-reflection. Our work builds on these observations but shifts the focus from \emph{whether} they fail to \emph{why} they fail through the lens of attribution.

\paragraph{Bias in LLM Evaluation} 
LLMs have been shown to exhibit "self-preference bias," where they prefer their own generated style or outputs in pairwise comparisons \cite{zheng2023judging}. Furthermore, \citet{source2025framing} demonstrated that source framing (e.g., attributing text to a specific nationality or model) triggers systematic shifts in evaluation scores. We extend this line of inquiry to the domain of logical reasoning and error detection.

\paragraph{Context and Anchoring}
The phenomenon of "sycophancy" in LLMs—the tendency to agree with the user's or their own previous stances—suggests that context plays a crucial role in evaluation quality. Works like \emph{Cross-Context Review} \cite{cross2026} emphasize that separating the generation and review sessions can mitigate these effects. We systematically test these session-level variables alongside identity labels.

\section{Experimental Setup}
\label{sec:setup}

\subsection{Identity Swap Critique Framework}
The core of our methodology involves a controlled manipulation of the critic's perception of authorship. Given a fixed problem $P$ and a flawed solution $S$ (generated by model $M_{solver}$), we prompt a critic model $M_{critic}$ to evaluate $S$ under five distinct conditions:
\begin{enumerate}
    \item \textbf{Self}: "This is your previous solution."
    \item \textbf{Other}: "This solution was generated by GPT-4."
    \item \textbf{Weak}: "This solution was generated by a smaller, 7B-parameter model."
    \item \textbf{Anonymous}: "Here is a solution for the given problem."
    \item \textbf{Human}: "A student wrote this solution."
\end{enumerate}

By keeping $S$ identical across all conditions, any variance in error detection accuracy can be directly attributed to the authorship framing.

\subsection{Datasets and Metrics}
We utilize three primary reasoning benchmarks: \textbf{GSM8K} for linguistic-math reasoning, \textbf{MATH} for complex step-by-step proofs, and \textbf{GPQA} for high-level scientific reasoning. 
Our evaluation metrics include:
\begin{itemize}
    \item \textbf{Detection Accuracy}: Success rate in identifying if an error exists.
    \item \textbf{Localization F1}: Accuracy in pinpointing the specific reasoning step where the error occurred.
    \item \textbf{Correction Rate}: The probability of successfully fixing the error given a correct detection.
\end{itemize}

[Table 1: Statistics of the datasets and error types used in experiments]

\section{Analysis: The Perspective Hypothesis}
\label{sec:analysis}

Our preliminary results suggest that models operate under a "trust-by-default" heuristic for their own outputs. When the model believes it is the author, it adopts a \emph{verification} mindset (looking for confirmation), whereas when it believes another entity is the author, it adopts a \emph{falsification} mindset (looking for flaws). 

This suggests that the "blind spot" is a result of cognitive alignment between the solver-state and the critic-state within the same model's parameter space. Separating these states—either by identity swapping or session isolation—is key to unlocking the model's latent reasoning potential.

...
\end{document}